\documentclass[11pt]{article}

\usepackage[utf8]{inputenc}
\usepackage[T1]{fontenc}
\usepackage[spanish,es-nodecimaldot]{babel}
\usepackage{amsmath,amssymb,mathtools}
\usepackage{geometry}
\usepackage{booktabs}
\usepackage{array}
\usepackage{enumitem}
\usepackage{hyperref}
\usepackage{xcolor}
\usepackage{siunitx}

\geometry{margin=2.5cm}
\hypersetup{
    colorlinks=true,
    linkcolor=blue!50!black,
    urlcolor=blue!50!black
}

\title{Derivaci\'on Mec\'anica Rigurosa\\Plataforma Antivibratoria Activa}
\author{Versi\'on corregida para estudio e informe}
\date{15 de abril de 2026}

\begin{document}

\maketitle
\tableofcontents
\newpage

\section{Punto de partida: qu\'e significa \textit{m\'as rigor} en este proyecto}

Antes de empezar hay que aclarar una idea importante: \textbf{no existe un modelo sin supuestos}. Todo modelo f\'isico hace supuestos. Lo correcto en un proyecto acad\'emico no es decir ``no hago supuestos'', sino:

\begin{itemize}[leftmargin=1.2cm]
\item hacer el \textbf{menor n\'umero posible} de supuestos;
\item que esos supuestos sean \textbf{expl\'icitos};
\item que cada supuesto tenga una \textbf{justificaci\'on f\'isica}.
\end{itemize}

En esta versi\'on voy a formular un modelo m\'as riguroso que el anterior, usando:

\begin{itemize}[leftmargin=1.2cm]
\item segunda ley de Newton para la parte mec\'anica;
\item ley de Hooke para el resorte;
\item ley del amortiguador viscoso;
\item ley de voltajes de Kirchhoff para un actuador electromec\'anico sencillo;
\item relaci\'on fuerza-corriente del actuador;
\item linealizaci\'on por Jacobianos alrededor del equilibrio.
\end{itemize}

Esta formulaci\'on tiene una ventaja grande: el sistema queda de tercer orden \textbf{sin inventar un estado de fuerza}, sino de manera directa a partir de una variable el\'ectrica interna del actuador.

\section{Qu\'e sistema f\'isico estamos modelando}

Se modela una \textbf{plataforma antivibratoria activa}. El sistema tiene:

\begin{itemize}[leftmargin=1.2cm]
\item una base inferior que se mueve seg\'un una perturbaci\'on externa;
\item una plataforma superior con masa concentrada;
\item un resorte entre base y plataforma;
\item un amortiguador entre base y plataforma;
\item un actuador lineal electromec\'anico que aplica fuerza entre base y plataforma.
\end{itemize}

Si quisi\'eramos una interpretaci\'on concreta, el actuador puede verse como un \textbf{voice coil}, un actuador lineal magn\'etico o un actuador equivalente cuyo esfuerzo es proporcional a la corriente.

\section{Variables y coordenadas}

\subsection{Coordenadas absolutas}

Primero definimos:

\begin{itemize}[leftmargin=1.2cm]
\item $x(t)$: posici\'on absoluta de la plataforma.
\item $z(t)$: posici\'on absoluta de la base.
\end{itemize}

\subsection{Coordenada relativa}

Para este tipo de sistema es muy conveniente definir la coordenada relativa:
\begin{equation}
q(t) = x(t) - z(t).
\label{eq:qdef}
\end{equation}

Esta coordenada tiene una interpretaci\'on f\'isica inmediata: es la deformaci\'on total de la suspensi\'on. De hecho:

\begin{itemize}[leftmargin=1.2cm]
\item si $q=0$, la plataforma y la base tienen la misma posici\'on;
\item si $q>0$, la plataforma est\'a desplazada por encima de la base;
\item si $q<0$, la plataforma est\'a desplazada por debajo de la base.
\end{itemize}

Derivando \eqref{eq:qdef}:
\begin{equation}
\dot{q}(t)=\dot{x}(t)-\dot{z}(t),
\label{eq:qdot}
\end{equation}
\begin{equation}
\ddot{q}(t)=\ddot{x}(t)-\ddot{z}(t).
\label{eq:qddot}
\end{equation}

\subsection{Por qu\'e conviene usar $q$}

Porque:

\begin{itemize}[leftmargin=1.2cm]
\item el resorte y el amortiguador dependen naturalmente de la deformaci\'on y velocidad relativas;
\item la ecuaci\'on mec\'anica queda m\'as limpia;
\item la perturbaci\'on de base entra de forma f\'isicamente clara.
\end{itemize}

\section{Derivaci\'on mec\'anica exacta del subsistema masa-resorte-amortiguador}

\subsection{Segunda ley de Newton}

Sobre la masa de la plataforma se aplica:
\begin{equation}
\sum F = m\ddot{x}.
\label{eq:newton}
\end{equation}

Las fuerzas que act\'uan sobre la plataforma son:

\begin{itemize}[leftmargin=1.2cm]
\item fuerza del resorte;
\item fuerza del amortiguador;
\item fuerza del actuador.
\end{itemize}

Por ahora voy a derivar la parte lineal m\'inima rigurosa. Despu\'es mostrar\'e c\'omo se agregan no linealidades realistas.

\subsection{Fuerza del resorte}

La ley de Hooke dice:
\begin{equation}
F_s = k\,\Delta \ell,
\end{equation}
donde $\Delta \ell$ es la deformaci\'on del resorte respecto a su longitud de equilibrio.

En nuestro sistema, la deformaci\'on del resorte es exactamente la diferencia de posiciones entre la plataforma y la base:
\[
\Delta \ell = x-z = q.
\]

Por tanto, la fuerza del resorte sobre la masa es restauradora:
\begin{equation}
F_s = -k(x-z) = -kq.
\label{eq:Fs}
\end{equation}

\subsection{Fuerza del amortiguador}

Para un amortiguador viscoso ideal:
\begin{equation}
F_d = c\,\Delta \dot{\ell},
\end{equation}
donde $\Delta \dot{\ell}$ es la velocidad relativa entre sus extremos.

Aqu\'i:
\[
\Delta \dot{\ell} = \dot{x}-\dot{z} = \dot{q}.
\]

La fuerza que se opone al movimiento relativo queda:
\begin{equation}
F_d = -c(\dot{x}-\dot{z}) = -c\dot{q}.
\label{eq:Fd}
\end{equation}

\subsection{Fuerza del actuador}

La fuerza del actuador se denota por:
\begin{equation}
F_a(t).
\end{equation}

La convenci\'on de signos que usaremos es: $F_a>0$ si el actuador empuja la plataforma en la direcci\'on positiva de $x$.

\subsection{Ecuaci\'on mec\'anica en coordenada absoluta}

Sustituyendo \eqref{eq:Fs} y \eqref{eq:Fd} en \eqref{eq:newton}:
\begin{equation}
m\ddot{x} = -k(x-z)-c(\dot{x}-\dot{z}) + F_a.
\label{eq:absolute_mech}
\end{equation}

\subsection{Pasar a coordenada relativa}

Usando \eqref{eq:qdef}--\eqref{eq:qddot}:

\[
x = q+z,\qquad
\dot{x}=\dot{q}+\dot{z},\qquad
\ddot{x}=\ddot{q}+\ddot{z}.
\]

Sustituyendo en \eqref{eq:absolute_mech}:
\begin{equation}
m(\ddot{q}+\ddot{z}) = -kq - c\dot{q} + F_a.
\end{equation}

Reordenando:
\begin{equation}
m\ddot{q} + c\dot{q} + kq = F_a - m\ddot{z}.
\label{eq:relative_mech}
\end{equation}

\subsection{Interpretaci\'on de \eqref{eq:relative_mech}}

Esta ecuaci\'on es muy importante porque muestra algo f\'isicamente elegante:

\begin{itemize}[leftmargin=1.2cm]
\item la variable interna del sistema es la deformaci\'on relativa $q$;
\item el resorte y el amortiguador aparecen de forma natural en $q$ y $\dot q$;
\item la perturbaci\'on de base entra como una \textbf{aceleraci\'on de base} $-\!m\ddot z$.
\end{itemize}

Es decir, desde la perspectiva relativa, la base excitada act\'ua como una fuerza inercial equivalente.

\section{Derivaci\'on del actuador desde leyes electromagn\'eticas sencillas}

Aqu\'i est\'a una de las partes que quer\'ias con m\'as rigor: \textbf{de d\'onde sale la ecuaci\'on del actuador}.

\subsection{Modelo m\'inimo de un actuador lineal electromec\'anico}

Para un actuador lineal de tipo electromagn\'etico, una ley bastante est\'andar es:
\begin{equation}
F_a = K_t i,
\label{eq:force_current}
\end{equation}
donde:

\begin{itemize}[leftmargin=1.2cm]
\item $i(t)$ es la corriente del actuador;
\item $K_t$ es la constante de fuerza del actuador.
\end{itemize}

Esto significa que la fuerza es proporcional a la corriente.

\subsection{Ley de voltajes de Kirchhoff}

En el circuito el\'ectrico del actuador se aplica KVL:
\begin{equation}
u = L\dot{i} + Ri + e_b,
\label{eq:kvl}
\end{equation}
donde:

\begin{itemize}[leftmargin=1.2cm]
\item $u(t)$ es el voltaje aplicado;
\item $L$ es la inductancia;
\item $R$ es la resistencia;
\item $e_b$ es la fuerza contraelectromotriz (back-EMF).
\end{itemize}

\subsection{Back-EMF}

Cuando el actuador se mueve, genera una fuerza contraelectromotriz proporcional a la velocidad relativa. Como el actuador trabaja entre la plataforma y la base, la velocidad relevante es la relativa:
\begin{equation}
e_b = K_e \dot{q},
\label{eq:bemf}
\end{equation}
donde $K_e$ es la constante de back-EMF.

\subsection{Ecuaci\'on el\'ectrica completa}

Sustituyendo \eqref{eq:bemf} en \eqref{eq:kvl}:
\begin{equation}
u = L\dot{i} + Ri + K_e\dot{q}.
\label{eq:electric}
\end{equation}

Reordenando:
\begin{equation}
L\dot{i} + Ri + K_e\dot{q} = u.
\label{eq:electric_reordered}
\end{equation}

\subsection{Interpretaci\'on f\'isica}

Esta ecuaci\'on dice:

\begin{itemize}[leftmargin=1.2cm]
\item parte del voltaje se usa en cambiar la corriente ($L\dot{i}$);
\item parte del voltaje cae en la resistencia ($Ri$);
\item parte del voltaje se ``pierde'' por el movimiento del actuador ($K_e\dot{q}$).
\end{itemize}

Eso es mucho m\'as riguroso que simplemente asumir un primer orden sin explicar de d\'onde sale.

\section{Modelo electrome\'canico acoplado}

Ahora unimos la parte mec\'anica \eqref{eq:relative_mech} con la parte el\'ectrica \eqref{eq:electric_reordered}.

\subsection{Sistema completo}

Usando \eqref{eq:force_current}, la ecuaci\'on mec\'anica queda:
\begin{equation}
m\ddot{q} + c\dot{q} + kq = K_t i - m\ddot{z}.
\label{eq:full_mech}
\end{equation}

Y la ecuaci\'on el\'ectrica es:
\begin{equation}
L\dot{i} + Ri + K_e\dot{q} = u.
\label{eq:full_elec}
\end{equation}

Este ya es un modelo \textbf{f\'isicamente riguroso, lineal y de tercer orden}.

\section{Estados del sistema}

Escogemos como estados:
\begin{equation}
x_1 = q,\qquad x_2 = \dot{q},\qquad x_3 = i.
\end{equation}

Entonces:
\begin{equation}
\dot{x}_1 = x_2,
\label{eq:ss1}
\end{equation}
\begin{equation}
\dot{x}_2 = -\frac{k}{m}x_1 - \frac{c}{m}x_2 + \frac{K_t}{m}x_3 - \ddot{z},
\label{eq:ss2}
\end{equation}
\begin{equation}
\dot{x}_3 = -\frac{K_e}{L}x_2 - \frac{R}{L}x_3 + \frac{1}{L}u.
\label{eq:ss3}
\end{equation}

En forma matricial:
\begin{equation}
\dot{x} =
\underbrace{
\begin{bmatrix}
0 & 1 & 0\\
-\dfrac{k}{m} & -\dfrac{c}{m} & \dfrac{K_t}{m}\\
0 & -\dfrac{K_e}{L} & -\dfrac{R}{L}
\end{bmatrix}}_{A}
x
+
\underbrace{
\begin{bmatrix}
0\\
0\\
\dfrac{1}{L}
\end{bmatrix}}_{B}
u
+
\underbrace{
\begin{bmatrix}
0\\
-1\\
0
\end{bmatrix}}_{E}
\ddot{z}.
\label{eq:ss_matrix}
\end{equation}

Si escogemos como salida la posici\'on relativa:
\[
y=q=x_1,
\]
entonces:
\begin{equation}
y=
\underbrace{
\begin{bmatrix}
1 & 0 & 0
\end{bmatrix}}_{C}
x.
\end{equation}

Si quieres la posici\'on absoluta de la plataforma, entonces:
\[
y=x=q+z.
\]

\section{Funci\'on de transferencia del modelo riguroso}

Para dise\~nar el controlador se toma primero la perturbaci\'on de base nula:
\[
z(t)=0.
\]

Entonces \eqref{eq:full_mech} y \eqref{eq:full_elec} quedan:
\begin{equation}
m\ddot{q} + c\dot{q} + kq = K_t i,
\label{eq:no_dist_mech}
\end{equation}
\begin{equation}
L\dot{i} + Ri + K_e\dot{q} = u.
\label{eq:no_dist_elec}
\end{equation}

Aplicando Laplace con condiciones iniciales nulas:
\[
(ms^2+cs+k)Q(s)=K_t I(s),
\]
\[
(Ls+R)I(s)+K_e s Q(s)=U(s).
\]

Despejando $I(s)$ de la primera:
\[
I(s)=\frac{ms^2+cs+k}{K_t}Q(s).
\]

Sustituyendo en la segunda:
\[
(Ls+R)\frac{ms^2+cs+k}{K_t}Q(s) + K_e s Q(s) = U(s).
\]

Multiplicando por $K_t$:
\[
\left[(Ls+R)(ms^2+cs+k)+K_tK_e s\right]Q(s)=K_tU(s).
\]

Por tanto:
\begin{equation}
G(s)=\frac{Q(s)}{U(s)}
=
\frac{K_t}
(Ls+R)(ms^2+cs+k)+K_tK_e s.
\label{eq:tf_exact}
\end{equation}

Expandiendo:
\begin{equation}
G(s)=\frac{K_t}
{Lm\,s^3 + (Lc+Rm)s^2 + (Lk+Rc+K_tK_e)s + Rk}.
\label{eq:tf_expanded}
\end{equation}

\subsection{Conclusi\'on}

La planta ya es \textbf{de tercer orden} y eso sale directamente de:

\begin{itemize}[leftmargin=1.2cm]
\item dos estados mec\'anicos ($q$ y $\dot q$);
\item un estado el\'ectrico ($i$).
\end{itemize}

\section{An\'alisis de estabilidad del modelo lineal riguroso}

La ecuaci\'on caracter\'istica es:
\[
Lm\,s^3 + (Lc+Rm)s^2 + (Lk+Rc+K_tK_e)s + Rk = 0.
\]

Para un c\'ubico
\[
a_3s^3+a_2s^2+a_1s+a_0,
\]
Routh-Hurwitz exige:

\begin{equation}
a_3>0,\quad a_2>0,\quad a_1>0,\quad a_0>0,\quad a_2a_1>a_3a_0.
\label{eq:rh}
\end{equation}

Aqu\'i:
\[
a_3=Lm,\qquad
a_2=Lc+Rm,\qquad
a_1=Lk+Rc+K_tK_e,\qquad
a_0=Rk.
\]

Si $L,m,c,R,k,K_t,K_e$ son positivos, entonces los cuatro coeficientes son positivos. La condici\'on no trivial es:
\begin{equation}
(Lc+Rm)(Lk+Rc+K_tK_e) > LmRk.
\label{eq:rh_condition}
\end{equation}

Esa es la condici\'on general de estabilidad del modelo lineal exacto.

\section{C\'omo agregar realismo sin perder rigor}

Hasta aqu\'i el modelo es ya muy riguroso y de tercer orden. Ahora bien, la actividad pide que el modelo realista incluya efectos adicionales que no est\'en en el modelo aproximado. Para eso se pueden a\~nadir:

\begin{itemize}[leftmargin=1.2cm]
\item rigidez c\'ubica;
\item fricci\'on seca regularizada;
\item saturaci\'on del actuador;
\item topes mec\'anicos.
\end{itemize}

Entonces el modelo realista puede escribirse como:
\begin{equation}
m\ddot{q} + c\dot{q} + kq + k_3q^3
+ F_c\tanh\left(\frac{\dot q}{v_\varepsilon}\right)
+ F_{\text{stop}}(q,\dot q)
= K_t i - m\ddot z,
\label{eq:nonlinear_mech}
\end{equation}
\begin{equation}
L\dot i + Ri + K_e\dot q = \mathrm{sat}(u).
\label{eq:nonlinear_elec}
\end{equation}

\subsection{Qu\'e parte es realmente ``mucho m\'as real''}

El modelo en Simscape deber\'ia representar precisamente \eqref{eq:nonlinear_mech}--\eqref{eq:nonlinear_elec} o una versi\'on equivalente de ellas usando bloques f\'isicos.

Es decir:

\begin{itemize}[leftmargin=1.2cm]
\item el \textbf{modelo realista} es el que simulas en Simscape;
\item el \textbf{modelo nominal} es el lineal, obtenido por linealizaci\'on y usado para dise\~no.
\end{itemize}

\section{Linealizaci\'on del modelo no lineal}

Ahora hacemos la parte que pediste con m\'as detalle: \textbf{c\'omo se linealiza exactamente}.

\subsection{Definir el vector de estados}

Tomamos:
\[
x=
\begin{bmatrix}
q\\
v\\
i
\end{bmatrix},
\qquad v=\dot q.
\]

Entonces el sistema no lineal queda:
\begin{equation}
\dot x_1 = x_2,
\end{equation}
\begin{equation}
\dot x_2 = \frac{1}{m}\left(
-kx_1 - cx_2 - k_3x_1^3 - F_c\tanh\left(\frac{x_2}{v_\varepsilon}\right)
-F_{\text{stop}}(x_1,x_2) + K_t x_3 - m\ddot z
\right),
\end{equation}
\begin{equation}
\dot x_3 = \frac{1}{L}\left(
-K_e x_2 - R x_3 + \mathrm{sat}(u)
\right).
\end{equation}

\subsection{Punto de equilibrio}

Tomamos el equilibrio:
\[
x^\star=
\begin{bmatrix}
0\\
0\\
0
\end{bmatrix},
\qquad
u^\star=0,
\qquad
\ddot z^\star=0.
\]

\subsection{Jacobiano}

La linealizaci\'on se obtiene con:
\[
A = \left.\frac{\partial f}{\partial x}\right|_{x^\star,u^\star},
\qquad
B = \left.\frac{\partial f}{\partial u}\right|_{x^\star,u^\star}.
\]

\subsection{Primera ecuaci\'on}

Como
\[
f_1(x,u)=x_2,
\]
sus derivadas son:
\[
\frac{\partial f_1}{\partial x_1}=0,\qquad
\frac{\partial f_1}{\partial x_2}=1,\qquad
\frac{\partial f_1}{\partial x_3}=0,\qquad
\frac{\partial f_1}{\partial u}=0.
\]

\subsection{Segunda ecuaci\'on}

Sea
\[
f_2(x,u)=
\frac{1}{m}\left(
-kx_1 - cx_2 - k_3x_1^3 - F_c\tanh\left(\frac{x_2}{v_\varepsilon}\right)
-F_{\text{stop}}(x_1,x_2) + K_t x_3 - m\ddot z
\right).
\]

Derivamos t\'ermino a t\'ermino.

\subsubsection{Derivada respecto a $x_1$}

\[
\frac{\partial}{\partial x_1}(-kx_1)=-k,
\qquad
\frac{\partial}{\partial x_1}(-k_3x_1^3)=-3k_3x_1^2.
\]

Evaluado en $x_1^\star=0$:
\[
-3k_3(x_1^\star)^2 = 0.
\]

Si el tope no est\'a activo en el equilibrio, entonces:
\[
\left.\frac{\partial F_{\text{stop}}}{\partial x_1}\right|_\star = 0.
\]

Por tanto:
\begin{equation}
\left.\frac{\partial f_2}{\partial x_1}\right|_\star = -\frac{k}{m}.
\end{equation}

\subsubsection{Derivada respecto a $x_2$}

\[
\frac{\partial}{\partial x_2}(-cx_2)=-c.
\]

Para el t\'ermino de fricci\'on regularizada:
\[
\frac{\partial}{\partial x_2}
\left[-F_c\tanh\left(\frac{x_2}{v_\varepsilon}\right)\right]
=
-F_c\frac{1}{v_\varepsilon}\operatorname{sech}^2\left(\frac{x_2}{v_\varepsilon}\right).
\]

Evaluado en $x_2^\star=0$:
\[
\operatorname{sech}^2(0)=1,
\]
entonces:
\[
\left.\frac{\partial}{\partial x_2}
\left[-F_c\tanh\left(\frac{x_2}{v_\varepsilon}\right)\right]\right|_\star
=
-\frac{F_c}{v_\varepsilon}.
\]

Si el tope no est\'a activo:
\[
\left.\frac{\partial F_{\text{stop}}}{\partial x_2}\right|_\star=0.
\]

Por tanto:
\begin{equation}
\left.\frac{\partial f_2}{\partial x_2}\right|_\star
=
-\frac{1}{m}\left(c+\frac{F_c}{v_\varepsilon}\right).
\end{equation}

\subsubsection{Derivada respecto a $x_3$}

\[
\left.\frac{\partial f_2}{\partial x_3}\right|_\star = \frac{K_t}{m}.
\]

\subsection{Tercera ecuaci\'on}

Sea
\[
f_3(x,u)=\frac{1}{L}\left(-K_e x_2 - Rx_3 + \mathrm{sat}(u)\right).
\]

Si linealizamos en una zona donde la saturaci\'on no est\'a activa, entonces localmente:
\[
\mathrm{sat}(u)\approx u.
\]

Por tanto:
\[
\left.\frac{\partial f_3}{\partial x_1}\right|_\star = 0,
\qquad
\left.\frac{\partial f_3}{\partial x_2}\right|_\star = -\frac{K_e}{L},
\qquad
\left.\frac{\partial f_3}{\partial x_3}\right|_\star = -\frac{R}{L},
\qquad
\left.\frac{\partial f_3}{\partial u}\right|_\star = \frac{1}{L}.
\]

\subsection{Matrices linealizadas}

Entonces:
\begin{equation}
A=
\begin{bmatrix}
0 & 1 & 0\\[4pt]
-\dfrac{k}{m} & -\dfrac{1}{m}\left(c+\dfrac{F_c}{v_\varepsilon}\right) & \dfrac{K_t}{m}\\[8pt]
0 & -\dfrac{K_e}{L} & -\dfrac{R}{L}
\end{bmatrix},
\qquad
B=
\begin{bmatrix}
0\\
0\\
\dfrac{1}{L}
\end{bmatrix}.
\label{eq:AB_lin}
\end{equation}

\subsection{Observaci\'on importante sobre la fricci\'on}

Si no quieres que el modelo nominal herede un t\'ermino extra de amortiguamiento por la regularizaci\'on $\tanh$, entonces puedes tomar la decisi\'on metodol\'ogica de:

\begin{itemize}[leftmargin=1.2cm]
\item incluir la fricci\'on solo en el modelo realista;
\item y omitirla del modelo nominal lineal.
\end{itemize}

Eso tambi\'en es v\'alido y, de hecho, puede ser m\'as limpio para dise\~no.

\section{C\'omo se obtiene el modelo simplificado que usamos antes}

Ahora conectamos este modelo riguroso con el modelo simplificado anterior.

\subsection{Hip\'otesis de simplificaci\'on}

Si en el rango de operaci\'on del dise\~no:

\begin{itemize}[leftmargin=1.2cm]
\item la velocidad relativa es peque\~na;
\item el back-EMF $K_e\dot q$ es mucho menor que $u$ y $Ri$;
\item y queremos trabajar con una variable interna fuerza en lugar de corriente,
\end{itemize}
entonces de \eqref{eq:electric_reordered} aproximamos:
\[
L\dot i + Ri \approx u.
\]

Como $F_a=K_t i$, entonces:
\[
\dot F_a = K_t \dot i.
\]

Multiplicando la ecuaci\'on aproximada por $K_t/R$:
\[
\frac{L}{R}\dot F_a + F_a = \frac{K_t}{R}u.
\]

Definiendo:
\[
\tau_a = \frac{L}{R},
\qquad
K_a = \frac{K_t}{R},
\]
obtenemos:
\begin{equation}
\tau_a\dot F_a + F_a = K_a u.
\label{eq:simplified_force_model}
\end{equation}

Esa es exactamente la forma simplificada que hab\'iamos usado antes.

\subsection{Interpretaci\'on}

Entonces el modelo anterior no estaba ``inventado'', sino que puede verse como una \textbf{aproximaci\'on de primer orden del subsistema el\'ectrico del actuador} cuando se desprecia el back-EMF.

\section{Respuesta directa a tu duda principal}

La secuencia correcta del proyecto es:

\begin{enumerate}[leftmargin=1.2cm]
\item plantear un modelo f\'isico del sistema;
\item derivar sus ecuaciones usando leyes f\'isicas;
\item construir un modelo realista de simulaci\'on;
\item linealizar alrededor del equilibrio para obtener el modelo nominal;
\item dise\~nar el controlador con ese modelo nominal;
\item comparar el modelo nominal y el modelo realista en simulaci\'on.
\end{enumerate}

O, dicho de forma muy clara:

\begin{quote}
\textbf{S\'i: primero planteamos el sistema realista, luego lo linealizamos para obtener el modelo nominal, y despu\'es lo comparamos con la versi\'on m\'as realista que simulamos en Simulink/Simscape.}
\end{quote}

La correcci\'on importante es que \textbf{Simscape no te ``entrega'' la ecuaci\'on de tercer orden}. Esa ecuaci\'on sale de tu derivaci\'on matem\'atica.

\section{Conclusi\'on}

La formulaci\'on m\'as rigurosa del proyecto, usando solo herramientas razonables del curso, es:

\begin{itemize}[leftmargin=1.2cm]
\item parte mec\'anica obtenida con Newton;
\item resorte obtenido con Hooke;
\item amortiguador con ley viscosa;
\item actuador obtenido con fuerza proporcional a corriente y KVL;
\item planta lineal exacta de tercer orden;
\item versi\'on realista obtenida agregando saturaci\'on, fricci\'on y topes;
\item linealizaci\'on por Jacobianos alrededor del equilibrio.
\end{itemize}

Si quieres m\'aximo rigor en el informe, esta es la versi\'on que yo usar\'ia como base te\'orica.

\end{document}
